\section{Theory Questions}

\subsection{a) Describe Timer Mode, PWM Mode, and Counter Mode}

\subsubsection*{Timer Mode}
Timer Mode uses an internal clock source, like the system clock, to increment a counter register at regular intervals. 
The Clock Signal passes through a Prescaler (divider), which can be configured to achieve the desired time increment. 
When the Counter reaches a predefined Threshold it generates an interrupt or sets a Flag, to notify the Program, that a specific time has passed.

This is typically used:
\begin{itemize}
    \item in delays or timeouts
    \item to generate periodic events (function calls, increments etc.)
    \item to time simple tasks such as LED Blinking
\end{itemize}

\subsubsection*{Counter Mode}
Here the Timer does not use the Internal Clock signal. Instead it counts external events or internal Triggers.
Each incoming event or Pulse increases the counter.

This is typically used:
\begin{itemize}
    \item to count Button presses
    \item to count rotations of a Motor or Servo
    \item to implement a Pulse Counter like in Task C in this exercise
\end{itemize}

\subsubsection*{PWM Mode}
In PWM (Pulse Width Modulation) Mode, the Timer produces a pulse-width modulation signal by comparing the counter value against a compare register.  
There are multiple different modes, like Fast-PWM, which starts at 0 and counts up to the maximum value and then resets, and Phase Correct PWM, which counts up and down for smoother signals.

\pagebreak

This is typically used:
\begin{itemize}
    \item to generate Analog Voltages in combination with a Filter
    \item to controll motor speed
    \item to dimm an LED by changing the Duty cycle
\end{itemize}

\subsection{b) Operation Mode: Timer with counter register}\label{task:counter_reg}
% Deine Antwort hier

In this Example the Timer is used in combination with the Counter Register.
As seen in the Block Schematic in Figure \ref{fig:timer_w_counter_reg} the Clock signal is divided and then fed into the Counter Register.

\begin{figure}[H]
    \centering
    \includegraphics[width=0.5\linewidth]{timer_w_counter_reg.png}
    \caption{Block Schematic of a Timer with counter Register}
    \label{fig:timer_w_counter_reg}
\end{figure}

As a numerical example we use:
\begin{align}
	f_{clk} &= \SI{10}{\kilo\hertz} \\
	D_{div} &= \num{250} \\
	f_{timer} &= \frac{f_{clk}}{D_{div}} 
	= \frac{\SI{10}{\kilo\hertz}}{\num{250}} 
	= \SI{40}{\hertz}
\end{align}

This means that the Prescaler reduces the Clock-tick timing from one tick every 
\begin{align}
\frac{1}{f_{clk}} = \SI{100}{\micro\second}
\end{align}

to one tick every 
\begin{align}
\frac{1}{f_{timer}} = \SI{25}{\milli\second}
\end{align}

The Counter value at $t = \SI{0}{\second}$ is \num{65}. The maximum value of the 16-bit counter is $2^{16} = \num{65536}$. This means that the overflow occurs after 
\begin{align}
2^{16} - 65 = \num{65471} \text{ counts}.
\end{align} 

Multiplying this by the time per count gives the time left until the overflow occurs:  
\begin{align}
\num{65471} \cdot \SI{25}{\milli\second} = \SI{1637}{\second} \approx \SI{27.28}{\minute}
\end{align}

So after approximately half an hour the Timer overflows. After the Overflow, the Counter resets to 0 and either the Overflow Flag is set, or a Timer overflow interrupt is executed. 

\subsection{c)	Operation Mode: Timer with capture event} \label{task:capture_event}

Here the timer works in the normal timer mode, as described above in \ref{task:counter_reg}. But additionally there is the Capture register. After each capture event (external or internal) happens, the current Counter Register value is copied into the Capture register and a capture flag is set. The counter continues running normally, but the captured register value acts as a timestamp for the event. Therefore, between two capture register entries a time difference can be measured. A Schematic of this process can be seen in Figure \ref{fig:timer_w_capture_event}.

\begin{figure}[H]
    \centering
    \includegraphics[width=0.5\linewidth]{timer_w_capture_event.png}
    \caption{Block Schematic of a Timer with capture Event}
    \label{fig:timer_w_capture_event}
\end{figure}

As a numerical example we use:
\begin{align}
	f_{clk} &= \SI{200}{\mega\hertz} \\
	D_{div} &= \num{500000} \\
	t_{event} &= \SI{2}{\second}
\end{align}
The Timer frequency can be calculated the same way as in task \ref{task:counter_reg} with 
\begin{align}
f_{timer} = \frac{f_{clk}}{D_{div}} = \SI{400}{\hertz}
\end{align}

Therefore each tick length is 
\begin{align}
\frac{1}{\SI{400}{\hertz}} = \SI{2.5}{\milli\second}.
\end{align} 

This means after \SI{2}{\second} the counter has counted to
\begin{align}
N = \frac{t}{T_{tick}} = \frac{\SI{2}{\second}}{\SI{2.5}{\milli\second}} = \num{800}.
\end{align}

This has not resulted in an overflow, because \num{800} $\ll 2^{16}$; therefore, the overflow flag has not been set.  
At $t = \SI{2}{\second}$ the Capture register has the same value as the Counter Register (\num{800}).

The board can then calculate the timespan of this event using:
\begin{align}
\Delta t = \text{capture value} \cdot T_{tick} = \SI{2}{\second}.
\end{align}

and more generally, the time span of an event by this formula:
\begin{align}
	\Delta t = \frac{\text{capture register value}\cdot N }{f_{clk}}
\end{align}

\subsection{d)	Operation Mode: Timer with compare register}

The Timer here works the same way as in task \ref{task:capture_event} but additionally a compare Register holds a threshold value. If this value is reached, a compare event is triggered. This compare event can be an Interrupt or an external Signal to toggle something outside of the chip. 

\begin{figure}[H]
    \centering
    \includegraphics[width=0.5\linewidth]{timer_w_compare_reg.png}
    \caption{Block Schematic of a Timer with compare Register}
    \label{fig:timer_w_compare_reg}
\end{figure}

As a numerical example we use:
\begin{align}
	f_{clk} &= \SI{20}{\kilo\hertz} \\
	D_{div} &= \num{10000} \\
	\text{Threshold} &= \num{10}
\end{align}

The Compare Register value is set to 10, this means after 10 clock cycles, a Compare event is executed. The Counter Register and Capture Register start at $0$.

First we calculate the Tick time again as seen in the previous two tasks.
\begin{align}
T_{tick} = \frac{1}{\frac{f_{clk}}{D_{div}}} = \SI{0.5}{\second}
\end{align}

In the following Table \ref{tab:compare_reg} we trace every Register value and the Actions of the compare register.

\begin{table}[h]
	\centering
	\begin{tabular}{| c | c | c | c | c |}
		\hline
		Time [s] & Capture reg. & Counter reg & Compare reg & Action \\
		\hline
		0   & 0 & 0  & 10 & none \\
		0.5 & 0 & 1  & 10 & none \\
		1   & 0 & 2  & 10 & none \\
		1.5 & 0 & 3  & 10 & none \\
		2   & 0 & 4  & 10 & none \\
		2.5 & 0 & 5  & 10 & none \\
		3   & 0 & 6  & 10 & none \\
		3.5 & 0 & 7  & 10 & none \\
		4   & 0 & 8  & 10 & none \\
		4.5 & 0 & 9  & 10 & none \\
		5   & 0 & 10 & 10 & Compare action \\
		\hline
	\end{tabular}
	\caption{Timer trace with Compare Register, Counter Register, Capture Register until an Action is triggered}
	\label{tab:compare_reg}
\end{table}
